\icode{E-EMA-SI-L2-260625-001}
如图甲、乙所示，在\( - d \le x\le d \)、\( - d\le y\le d\)区域中的阴影部分，存在垂直\(xOy\)平面向里的匀强磁场，边长均为\(2d\)正方形线框与磁场边界重合。两线框都绕\(y\)轴以相同的角速度匀速转运，则甲、乙线框中产生的感应电动势最大值之比为\paren[]

\begin{tikzpicture}[font=\scriptsize,>=latex]
    \begin{scope}
        \draw[fill=gray!50,draw=white] (-1.2,1.2) rectangle(-.6,-1.2) (.6,1.2)  rectangle(1.2,-1.2);
        \foreach \x in{-.9,.9}
            \foreach \y in{-.8,-.4,.4,.8}
                \node[teal] at(\x,\y){\( \times \)};
        \draw[dashed] (-.6,1.2)--(-.6,-1.2)node[below]{\( -\frac{d}{2}\)} (.6,1.2)--(.6,-1.2)node[below]{\(\frac{d}{2}\)};
        \draw[thick] (-1.2,1.2) rectangle(1.2,-1.2);\draw[blue,->] (-1.6,0)--(1.6,0)node[right]{\(x\)};
        \draw[blue,->] (0,-1.8)--(0,2)node[left]{\(y\)};
        \draw[orange,->] (-.12,1.5) arc(130:430:.2 and .1);
        \node[blue] at(-.2,-.2){\(O\)};
        \node[below] at(0,-2){甲};
    \end{scope}
    \begin{scope}[xshift=4cm]
        \draw[fill=gray!50,draw=white](0,1.2)  rectangle(1.2,-1.2);
        \foreach \x in{.3,.9}
            \foreach \y in{-.8,-.4,.4,.8}
                \node[teal] at(\x,\y){\( \times \)};
        \draw[thick] (-1.2,1.2) rectangle(1.2,-1.2);\draw[blue,->] (-1.6,0)--(1.6,0)node[right]{\(x\)};
        \draw[blue,->] (0,-1.8)--(0,2)node[left]{\(y\)};
        \draw[orange,->] (-.12,1.5) arc(130:430:.2 and .1);
        \node[blue] at(-.2,-.2){\(O\)};
        \node[below] at(0,-2){乙};
    \end{scope}
\end{tikzpicture}
\begin{choices}
    \item \(1:1\)
    \item \(\sqrt{3}:2\)
    \item \(\sqrt{3}:1\)
    \item \(2:1\)
\end{choices}
\begin{answer}
    C\quad 提示：
\end{answer}